\documentclass[11pt]{article}
\usepackage[margin=1.1in]{geometry}
\usepackage{amsmath, amssymb, amsthm, mathtools}
\usepackage[colorlinks=true, linkcolor=blue, citecolor=blue]{hyperref}

\newtheorem{theorem}{Theorem}
\newtheorem{lemma}{Lemma}
\newtheorem{proposition}{Proposition}
\newtheorem{corollary}{Corollary}
\newtheorem{assumption}{Assumption}
\theoremstyle{definition}
\newtheorem{definition}{Definition}
\newtheorem{remark}{Remark}
\newtheorem{prediction}{Prediction}

\newcommand{\R}{\mathbb{R}}
\newcommand{\E}{\mathbb{E}}
\newcommand{\cA}{\mathcal{A}}
\newcommand{\cB}{\mathcal{B}}
\newcommand{\cP}{\mathcal{P}}
\newcommand{\cM}{\mathcal{M}}
\newcommand{\ip}[2]{\left\langle #1, #2 \right\rangle}
\newcommand{\KL}{\mathrm{KL}}
\DeclareMathOperator*{\argmax}{arg\,max}

\title{Why Magnetic Mirror Descent on Gaussian--Mixture Policies\\
Converges to Mixed Nash Equilibria of Continuous Games:\\
Lifted Monotonicity and Local Monotone Islands}
\author{Theory notes for \texttt{cont\_tinkering}}
\date{July 2026}

\begin{document}
\maketitle

\begin{abstract}
Magnetic Mirror Descent (MMD) provably converges to Nash equilibria only for
\emph{monotone} games, yet empirically it converges on continuous-action
zero-sum games with mixed equilibria when the policy is a $K$-component
Gaussian mixture (categorical head over component means/scales) --- a setting in
which the game, viewed in policy \emph{parameters}, is far from monotone. We
explain this with a three-layer argument. (1)~The mixed extension of \emph{any}
zero-sum game over probability measures is bilinear, hence its pseudo-gradient
operator is monotone --- in fact exactly skew; the measure-space MMD dynamics are
therefore squarely inside the standard convergence theory, and all
non-monotonicity is an artifact of the parameterization. (2)~The implemented
update is mirror descent in the Kullback--Leibler geometry of the hierarchical
policy (exact simplex MMD on the categorical head, Fisher--Rao natural gradient
on the Gaussian heads): it is the restriction of the lifted monotone flow to
the mixture submanifold. (3)~For zero-sum games the symmetrized Jacobian of the
parametric pseudo-gradient field is block-diagonal (cross-player blocks are
skew), so the field is monotone on any region where each player's payoff is
concave in its own parameters. At a Nash-realizing configuration each player's
\emph{effective landscape} --- its own payoff term plus the coupling evaluated
at the opponent's generalized moments --- has a strict local maximum at every
component mean, so the payoff is locally strongly concave in own parameters,
the field is locally (strongly, once the magnet is added) monotone, and MMD
converges linearly. Convergence is therefore \emph{local by nature}: other
locally-monotone ``islands'' (a mass-dominated decoy, symmetric ties,
simplex-boundary freezes) coexist, and the initialization selects among them.
Nothing in the development assumes the game is symmetric under exchange of the
players: the two players may have different action sets, different landscapes,
different feature maps, different numbers of mixture components and different
hyperparameters. Every quantitative prediction of this theory is verified
numerically in the accompanying experiments (\texttt{theory/exp1--exp5}).
\end{abstract}

\section{Setting and notation}

\subsection{Game class}\label{sec:class}

Throughout, player~$0$ is the maximizer with compact action set
$\cA \subset \R^{n_0}$ and player~$1$ is the minimizer with compact action set
$\cB \subset \R^{n_1}$. The two sets need not coincide, and $n_0 \ne n_1$ is
allowed. We consider three nested levels of structure, and we will be explicit
about which level each statement needs.

\paragraph{(G0) Arbitrary zero-sum.} $u : \cA \times \cB \to \R$ bounded and
measurable. Nothing else. This is all that Section~\ref{sec:layer1} uses.

\paragraph{(G1) Finite rank plus own landscapes.} There are $C^2$ functions
$\varphi_1,\dots,\varphi_m : \cA \to \R$ and $\psi_1,\dots,\psi_m : \cB \to \R$
and $C^2$ \emph{landscapes} $D_0 : \cA \to \R$, $D_1 : \cB \to \R$ with
\begin{equation}\label{eq:game}
  u(a,b) \;=\; D_0(a) \;-\; D_1(b) \;+\; \sum_{j=1}^{m}
  \varphi_j(a)\, \psi_j(b)
  \;=\; D_0(a) - D_1(b) + \ip{\varphi(a)}{\psi(b)} .
\end{equation}
The defining property is that the opponent enters each player's utility only
through the finite-dimensional \emph{moment vector}
$\E_\nu \psi \in \R^m$ (resp.\ $\E_\mu \varphi \in \R^m$); the rank $m$ may be
arbitrary but is finite. A coupling of the form $c\,\ip{f(a)}{M g(b)}$ with a
constant matrix $M$ and a scalar strength $c$ is the special case
$\varphi = c f$, $\psi = M g$, so writing the strength explicitly is a matter
of taste; we keep $c$ visible where the experiments sweep it, i.e.\ we write
$\varphi = c\, f$ with $\|f\|$ normalized. Section~\ref{sec:local} uses (G1)
only to make its hypotheses \emph{checkable}; the theorem itself is stated
under abstract hypotheses that (G1) implies.

\paragraph{(G2) The codebase class.} (G1) with $\cA = \cB = [a_{\min},
a_{\max}] \subset \R$, each landscape a finite sum of Gaussian bumps
$D_i(x) = \sum_j h^{(i)}_j \exp\big({-}(x - p^{(i)}_j)^2 / 2 (w^{(i)}_j)^2\big)$,
and $\varphi_j(a) = c\,(\iota(a)^j - t_j)$, $\psi_j(b) = \iota(b)^j - t'_j$
normalized monomials in a rescaled coordinate $\iota$, with the shifts
$t, t'$ chosen so that the equilibrium moments vanish (see
Remark~\ref{rem:centering}: the shifts are a normalization, not a restriction).
This class contains \texttt{MultiPointGame} and \texttt{DecoyWellGame}. It is
used only to obtain closed forms: the smoothed landscape, the smoothed
features, and the explicit constants of Lemma~\ref{lem:bias} and
Assumption~\ref{ass:curv}. The multidimensional separable case is
Remark~\ref{rem:multidim}.

\begin{remark}[No symmetry is assumed]\label{rem:asym}
The classical instances in the codebase happen to satisfy $D_0 = D_1$,
$\varphi = c\,f$, $\psi = f$, $\cA = \cB$, hence $u(a,b) = -u(b,a)$: the game
is symmetric under exchange of the players. \emph{None of the results below use
this.} What is used is that the game is \emph{zero-sum}: a single payoff $u$
enters the two players' objectives with opposite signs. Where the symmetric
case would let one write ``and player~1 is symmetric'', we instead carry the
second player's own objects ($D_1$, $\psi$, $K_1$, $\eta_1$, $\tau_1$, floor
$s_1$) explicitly. The one place where player-exchange symmetry does more than
save notation is discussed in Remark~\ref{rem:where-symmetry-helps}, and it
concerns the \emph{global} transit analysis, not any statement proved here.
\end{remark}

\begin{remark}[Centering convention]\label{rem:centering}
Under (G1), let $(\mu^\star, \nu^\star)$ be a mixed equilibrium and put
$\bar\psi \coloneqq \E_{\nu^\star} \psi$, $\bar\varphi \coloneqq
\E_{\mu^\star} \varphi$. Replacing
\[
  D_0 \leftarrow D_0 + \ip{\varphi(\cdot)}{\bar\psi}, \quad
  D_1 \leftarrow D_1 + \ip{\bar\varphi}{\psi(\cdot)}, \quad
  \varphi \leftarrow \varphi - \bar\varphi, \quad
  \psi \leftarrow \psi - \bar\psi
\]
changes $u$ by the constant $-\ip{\bar\varphi}{\bar\psi}$ and leaves the game
strategically identical, while enforcing
$\E_{\nu^\star}\psi = 0$ and $\E_{\mu^\star}\varphi = 0$. Hence the statement
``at the equilibrium the coupling term vanishes identically'' is a
\emph{normalization}, available in every game of class (G1), and not an extra
hypothesis --- in particular it is not a consequence of symmetry. What is
\emph{not} free, and is the real hypothesis of Section~\ref{sec:local}, is that
the atoms of $\mu^\star$ are strict local maxima of the \emph{recentered}
landscape $D_0$ (and the atoms of $\nu^\star$ strict local minima of the
recentered $D_1$). In the codebase this recentering is what the target shifts
$t_j$ implement.
\end{remark}

\begin{remark}[Scope map: how much structure each result needs]\label{rem:scope}
\begin{center}
\begin{tabular}{lll}
\hline
result & needs & does \emph{not} need\\
\hline
Lemma~\ref{lem:lift} (lifted skewness) & (G0), zero-sum & structure of $u$,
symmetry, $\cA = \cB$\\
Cor.~\ref{cor:measure-mmd} & (G0) + Lipschitz field & anything parametric\\
Lemma~\ref{lem:skewblock} (block-skew) & zero-sum, $C^2$ in parameters &
(G1), equal parameter dimensions\\
Lemma~\ref{lem:bias} (bias) & (G2) for the constant; (G1) + $C^2$ for
$O(s^2)$ & symmetry, equal floors\\
Thm.~\ref{thm:local} (islands) & Assumptions \ref{ass:reg}--\ref{ass:curv} &
(G1)/(G2); these are implied by, not equivalent to, (G1)\\
Prop.~\ref{prop:islands} (coexisting islands) & Assumptions
\ref{ass:reg}--\ref{ass:curv} & moment matching, symmetry\\
\hline
\end{tabular}
\end{center}
In words: Layers 1--2 are structure-free consequences of the zero-sum property;
Layer~3 needs only that each player's \emph{effective landscape} is $C^2$ and
strongly concave (resp.\ convex) near the realized atoms, uniformly in the
opponent's policy; class (G1) is what makes ``uniformly in the opponent's
policy'' checkable, because the opponent enters through $m$ numbers; class (G2)
is what makes the constants explicit.
\end{remark}

\subsection{Mixed extension}
For $\mu \in \cP(\cA)$, $\nu \in \cP(\cB)$ (Borel probability measures),
\begin{equation}\label{eq:mixed}
  U(\mu, \nu) = \E_{a \sim \mu,\, b \sim \nu}\, u(a,b) .
\end{equation}
By Glicksberg's theorem a mixed Nash equilibrium exists whenever $u$ is
continuous on the compact $\cA \times \cB$. We write the two \emph{effective
landscapes} --- the object each player actually climbs, given the opponent ---
as
\begin{equation}\label{eq:effective}
  Q_\nu(a) \coloneqq \int u(a,b)\, d\nu(b), \qquad
  W_\mu(b) \coloneqq -\int u(a,b)\, d\mu(a),
\end{equation}
so that player~0 maximizes $\E_\mu Q_\nu$ and player~1 maximizes
$\E_\nu W_\mu$. Under (G1) and up to additive constants (which are irrelevant
to both the dynamics and the equilibrium),
\begin{equation}\label{eq:effective-G1}
  Q_\nu(a) = D_0(a) + \ip{\varphi(a)}{\E_\nu \psi}, \qquad
  W_\mu(b) = D_1(b) - \ip{\E_\mu \varphi}{\psi(b)} .
\end{equation}
Note the \emph{opposite signs} of the two coupling terms: this is the point at
which a symmetric treatment silently conflates the two players. All statements
below are made separately for $Q$ and $W$.

\subsection{Policy class and the implemented update}
Player $i$ uses a $K_i$-component Gaussian mixture (the numbers $K_0, K_1$ need
not agree)
\begin{equation}
  \pi_{\theta_i} = \sum_{k=1}^{K_i} w^{(i)}_k\,
  \mathcal{N}\big(\mu^{(i)}_k, (\sigma^{(i)}_k)^2\big), \qquad
  \theta_i = (\ell^{(i)}, \mu^{(i)}, \rho^{(i)}), \quad
  w^{(i)} = \mathrm{softmax}(\ell^{(i)}),
\end{equation}
with $\sigma^{(i)}_k = e^{\rho^{(i)}_k}$ subject to a hard floor
$\sigma^{(i)}_k \ge s_i$ ($s_i = 10^{-3}$ in the code). The update
(\texttt{idealized\_mmd.py}) is, per player $i$, with step $\eta_i$, magnet
temperature $\tau_i$, magnet policy $\pi_{\bar\theta_i}$ refreshed every $T$
steps:
\begin{align}
  \text{categorical head:}\quad
  & \ell^{(i),+} = \frac{\eta_i\, q^{(i)} + \eta_i \tau_i \log \bar w^{(i)}
    + \log w^{(i)}}{1 + \eta_i\tau_i},
  \label{eq:cat}\\[2pt]
  & \qquad q^{(0)}_k = \E_{a \sim \mathcal{N}(\mu^{(0)}_k, (\sigma^{(0)}_k)^2)}
    Q_\nu(a), \qquad
    q^{(1)}_k = \E_{b \sim \mathcal{N}(\mu^{(1)}_k, (\sigma^{(1)}_k)^2)}
    W_\mu(b), \notag\\
  \text{Gaussian heads:}\quad
  & \mu^{(i),+} = \mu^{(i)} + \eta_i\, (\sigma^{(i)})^2 \odot
    \nabla_{\mu^{(i)}} \Phi_i, \qquad
    \rho^{(i),+} = \rho^{(i)} + \tfrac{\eta_i}{2}\,
    \nabla_{\rho^{(i)}} \Phi_i,
  \label{eq:gauss}
\end{align}
where $\Phi_0 = U - \tau_0 \KL(\pi_{\theta_0} \| \pi_{\bar\theta_0})$ and
$\Phi_1 = -U - \tau_1 \KL(\pi_{\theta_1} \| \pi_{\bar\theta_1})$ (each player
ascends its own objective), and the preconditioners $(\sigma^{(i)})^2$ and
$\tfrac12$ are the inverse Fisher metric of $\mathcal{N}(\mu, e^{\rho})$.
Equation~\eqref{eq:cat} is the \emph{exact closed-form} MMD proximal step on
the simplex (Sokota et al., 2023); equation~\eqref{eq:gauss} is its small-step
(Fisher--Rao natural gradient) limit for the location--scale heads. Thus each
player's update is mirror descent in the KL geometry of its own hierarchical
policy $(k \sim \mathrm{Cat}(w^{(i)}),\ x \sim \mathcal{N}(\mu^{(i)}_k,
(\sigma^{(i)}_k)^2))$. Note that \eqref{eq:cat}--\eqref{eq:gauss} are
\emph{decoupled across players} given the opponent's policy: nothing in the
algorithm presumes the two players run the same $\eta$, $\tau$, $K$ or floor.

\section{Layer 1: the lifted game is monotone --- exactly skew}
\label{sec:layer1}

\begin{definition}[Monotone game]
Let $z = (\mu, \nu) \in \cP(\cA) \times \cP(\cB)$ and let
$F(z) = \big({-}\nabla_{\mu} U(\mu,\nu),\ \nabla_{\nu} U(\mu,\nu)\big)$
be the descent field of the min--max problem $\max_\mu \min_\nu U$. The game is
\emph{monotone} on a convex set $S$ if
$\ip{F(z) - F(z')}{z - z'} \ge 0$ for all $z, z' \in S$.
\end{definition}

\begin{lemma}[Lifted monotonicity]\label{lem:lift}
For \emph{any} bounded measurable $u : \cA \times \cB \to \R$ --- in particular
with $\cA \ne \cB$, of different dimensions, and with no symmetry between the
players --- the mixed extension \eqref{eq:mixed} is affine in $\mu$ for fixed
$\nu$ and affine in $\nu$ for fixed $\mu$, and its pseudo-gradient field is
monotone on $\cP(\cA) \times \cP(\cB)$ with equality:
\[
  \ip{F(z) - F(z')}{z - z'} = 0 \qquad
  \forall\, z, z' \in \cP(\cA) \times \cP(\cB) .
\]
\end{lemma}

\begin{proof}
We spell the argument out in five steps. Throughout, the pairing between a
bounded measurable function $g$ and a finite signed measure $\delta$ is
$\ip{g}{\delta} \coloneqq \int g\, d\delta$ --- the measure plays the role of
the vector, the function the role of the gradient (in the finite-dimensional
discretization this is literally $\ip{g}{\delta} = \sum_i g_i \delta_i$, and
the whole proof below reduces to: $U(p,q) = p^\top\! A q$ with $A$ an arbitrary
\emph{rectangular} matrix, $F(p,q) = (-Aq,\, A^\top\! p)$, and
$-\delta p^\top\! A\, \delta q + \delta q^\top\! A^\top \delta p = 0$ because a
scalar equals its own transpose --- squareness of $A$, i.e.\ $|\cA| = |\cB|$,
is never used, and neither is $A = -A^\top$, i.e.\ symmetry of the game).

\emph{Step 1 (linearity and gradients).} By Fubini,
\[
  U(\mu, \nu) = \iint u(a,b)\, d\mu(a)\, d\nu(b)
  = \int \Big[\underbrace{\textstyle\int u(a,b)\, d\nu(b)}_{=\, Q_\nu(a)}
  \Big]\, d\mu(a) = \ip{Q_\nu}{\mu}.
\]
For fixed $\nu$, $U$ is therefore a \emph{linear} functional of $\mu$ with
coefficient function $Q_\nu$; its G\^ateaux derivative at any $\mu$ in the
direction of a signed measure $\delta$ is
$\tfrac{d}{dt} U(\mu + t\delta, \nu)\big|_{t=0} = \ip{Q_\nu}{\delta}$,
so $\nabla_\mu U = Q_\nu$ --- a function on $\cA$ only, independent of $\mu$
itself. In the same way $\nabla_\nu U = -W_\mu$, a function on $\cB$ only,
with $W_\mu$ as in \eqref{eq:effective}. The two objects live on different
spaces and are computed from different sections of $u$; no relation between
them is needed.

\emph{Step 2 (the field).} Player $0$ maximizes and player $1$ minimizes, so
the descent field of the min--max problem is
$F(\mu, \nu) = \big({-}Q_\nu,\ -W_\mu\big)$; note each block depends only on
the \emph{opponent's} measure.

\emph{Step 3 (blockwise expansion).} With $z = (\mu,\nu)$, $z' = (\mu',\nu')$,
$\delta\mu \coloneqq \mu - \mu'$, $\delta\nu \coloneqq \nu - \nu'$, the inner
product on the product space is the sum of the per-player pairings:
\[
  \ip{F(z) - F(z')}{z - z'}
  = \ip{-Q_\nu + Q_{\nu'}}{\delta\mu} + \ip{-W_\mu + W_{\mu'}}{\delta\nu}.
\]

\emph{Step 4 (each block is the same double integral).} By linearity of the
integral in the measure,
\[
  Q_\nu(a) - Q_{\nu'}(a) = \int u(a,b)\, d(\delta\nu)(b),
\]
and pairing this function of $a$ against $\delta\mu$ means integrating it
against $d(\delta\mu)(a)$:
\[
  \ip{Q_\nu - Q_{\nu'}}{\delta\mu}
  = \int\!\Big[\int u(a,b)\, d(\delta\nu)(b)\Big]\, d(\delta\mu)(a)
  = \iint u\; d(\delta\mu)\, d(\delta\nu) \eqqcolon I .
\]
The second block is the mirrored computation: $W_{\mu'}(b) - W_\mu(b) =
\int u(a,b)\, d(\delta\mu)(a)$, and pairing against $\delta\nu$ gives, by
Fubini, the \emph{identical} scalar $I$.

\emph{Step 5 (cancellation).} The first block carries the minus sign from the
field:
\[
  \ip{F(z) - F(z')}{z - z'} = -I + I = 0. \qedhere
\]
\end{proof}

\begin{corollary}\label{cor:measure-mmd}
Measure-space MMD (equivalently: tabular MMD on any finite discretization of
$\cA$ and $\cB$, with the two grids independent) on any game of class (G0)
satisfies the monotonicity assumption of Sokota et al.\ (2023, Thm.~3.4): with
entropy regularization $\tau > 0$ it converges linearly to the
$\tau$-regularized equilibrium provided the step size obeys the theorem's bound
(of the form $\eta \le \tau / L^2$ with $L$ the operator's Lipschitz constant),
and the magnet outer loop drives the regularized equilibria to a Nash
equilibrium.
\end{corollary}

\begin{remark}
Corollary~\ref{cor:measure-mmd} already dissolves half the puzzle: the games
here (including the ``non-monotone'' \textsc{Forsaken} saddle and the decoy
well) are \emph{monotone in measure space}. Classical non-convergence results
for these games concern \emph{pure-strategy} (point) dynamics. Experiment~1
confirms both halves: tabular MMD converges globally on all of them (even
initialized on the decoy --- slowly, gaining $\Theta(\Delta q / \tau)$ log-mass
per magnet cycle, but inevitably, because in measure space support never
dies), and violating the step-size bound makes it cycle even in the exact
lifted space, cleanly separating step-size failures from parameterization
failures.
\end{remark}

\section{Layer 2: zero-sum structure of the parametric field}

Restrict $U$ to the mixture families:
$\widehat{U}(\theta_0, \theta_1) \coloneqq U(\pi_{\theta_0}, \pi_{\theta_1})$
with $\theta_0 \in \R^{d_0}$, $\theta_1 \in \R^{d_1}$ and, in general,
$d_0 \ne d_1$ (different $K_i$, different action dimensions). Let
\[
  F(\theta_0,\theta_1) =
  \big({-}\nabla_{\theta_0}\widehat U,\ \nabla_{\theta_1}\widehat U\big)
\]
be the parametric descent field ($C^2$ on the interior of the parameter box).

\begin{lemma}[Block-skew structure]\label{lem:skewblock}
The Jacobian of $F$,
\[
  J = \begin{pmatrix}
    -\nabla^2_{\theta_0\theta_0} \widehat U & -\nabla^2_{\theta_0\theta_1} \widehat U\\[2pt]
    \ \ \nabla^2_{\theta_1\theta_0} \widehat U & \ \ \nabla^2_{\theta_1\theta_1} \widehat U
  \end{pmatrix}
  \in \R^{(d_0 + d_1) \times (d_0 + d_1)},
\]
has symmetric part
$\tfrac12 (J + J^\top) = \mathrm{blockdiag}\!\big({-}\nabla^2_{\theta_0\theta_0}
\widehat U,\ \nabla^2_{\theta_1\theta_1} \widehat U\big)$: the cross-player
blocks cancel exactly. Consequently $F$ is monotone on a convex set $S$ if and
only if $\widehat U(\cdot, \theta_1)$ is concave in $\theta_0$ and $\widehat
U(\theta_0, \cdot)$ is convex in $\theta_1$ throughout $S$; and $F$ is
$\alpha$-strongly monotone on $S$ iff the two own-blocks are uniformly
$\alpha$-definite there.
\end{lemma}

\begin{proof}
The off-diagonal blocks of $J$ are the $d_0 \times d_1$ matrix
$-\nabla^2_{\theta_0\theta_1}\widehat U$ and the $d_1 \times d_0$ matrix
$+\nabla^2_{\theta_1\theta_0}\widehat U = +\big(\nabla^2_{\theta_0\theta_1}
\widehat U\big)^{\!\top}$ (equality of mixed partials of the single scalar
$\widehat U$); in $\tfrac12(J + J^\top)$ they appear as
$\tfrac12\big({-}\nabla^2_{\theta_0\theta_1}\widehat U +
\nabla^2_{\theta_0\theta_1}\widehat U\big) = 0$. Monotonicity of a $C^1$ field
on a convex set is equivalent to positive semidefiniteness of the symmetrized
Jacobian on that set, which is now equivalent to the stated
concavity/convexity. \emph{This uses only that the game is zero-sum} --- the
single payoff $\widehat U$ appears with opposite signs for the two players, so
whatever cross-curvature exists is antisymmetric between the blocks. It does
not use $d_0 = d_1$, any correspondence between the two parameterizations, or
any symmetry of $u$; the cancellation is between a matrix and the transpose of
the \emph{same} matrix, which exists for rectangular blocks just as well.
\end{proof}

\begin{remark}[Where monotonicity is lost, and where it is not]
Lemma~\ref{lem:skewblock} localizes the entire failure of monotonicity in the
own-player curvature. The interaction between the players --- however strong the
coupling --- never breaks monotonicity by itself. The mixture parameterization
breaks it exactly where $\theta_i \mapsto \widehat U$ fails the required
own-curvature sign, i.e.\ where a component mean sits in a convex region of its
effective landscape \eqref{eq:effective}. Note also that the categorical heads
never contribute: $\widehat U$ is \emph{exactly linear} in each player's own
mixture weights $w^{(i)}$, and update~\eqref{eq:cat} operates on the simplex
where this linearity lives (the softmax non-concavity in $\ell$ is absorbed by
the mirror map). Because the lemma treats the two players independently, one
player may be in a monotone region while the other is not; the field is
monotone precisely on the intersection.
\end{remark}

\section{Layer 3: local monotone islands}\label{sec:local}

\subsection{Effective landscapes and Nash-realizing configurations}

Write $\Lambda_s \coloneqq \Lambda * \mathcal{N}(0, s^2)$ for the Gaussian
smoothing of a function $\Lambda$ at scale $s$; under (G2) both $D_{i,s}$ and
the smoothed features are available in closed form, which is what the code
computes. Fix an equilibrium $(\mu^\star, \nu^\star)$, recentered as in
Remark~\ref{rem:centering}, with finite supports
\[
  \mathrm{supp}(\mu^\star) = \{p_1, \dots, p_{K_0^\star}\} \subset \cA,
  \qquad
  \mathrm{supp}(\nu^\star) = \{r_1, \dots, r_{K_1^\star}\} \subset \cB,
\]
and weights $w^{(0)\star}, w^{(1)\star}$. The two supports are unrelated: they
live in different spaces and may have different cardinalities.

\begin{definition}[Nash-realizing configuration]
Given scales $s_0, s_1$, the configuration $z^\star_s = (\theta^\star_0,
\theta^\star_1)$ has $K_i = K_i^\star$, $\mathrm{softmax}(\ell^{(i)\star}) =
w^{(i)\star}$, one component mean of player~0 at each $p_k$, one component mean
of player~1 at each $r_l$, and $\sigma^{(i)} = s_i$ at the respective floor.
\end{definition}

\begin{assumption}[Regularity and identifiability]\label{ass:reg}
(i) $u$ is $C^2$ on $\cA \times \cB$, with $\sup \|\nabla^2 u\| < \infty$.
(ii) The $K_1^\star \times K_0^\star$ payoff matrix $A_{lk} = u(p_k, r_l)$ has
$(w^{(0)\star}, w^{(1)\star})$ as its \emph{unique} Nash equilibrium, and both
are fully mixed. (iii) Each $p_k$ is a strict local maximum of the recentered
$D_0$, each $r_l$ a strict local maximum of the recentered $D_1$.
\end{assumption}

Part (ii) is the general form of the Vandermonde argument in the
\texttt{MultiPointGame} docstrings: there, $A$ is generated by the moment
features and its nonsingularity follows from distinctness of the support
points. Stated as a property of the matrix $A$ it needs neither the monomial
structure nor $K_0^\star = K_1^\star$: it is the standard condition that the
support-restricted matrix game has a unique fully-mixed equilibrium.

\begin{lemma}[Representation and smoothing bias]\label{lem:bias}
Under Assumption~\ref{ass:reg} and class (G1), at $z^\star_s$ with
$K_i = K_i^\star$:
$\E_{\pi_{\theta_0^\star}} \varphi = O(s_0^2)$ and
$\E_{\pi_{\theta_1^\star}} \psi = O(s_1^2)$ (exactly $0$ for features affine in
the action), and the exploitability of the pair splits into a sum of
\emph{per-player} smoothing deficits:
\[
  \mathrm{expl}(z^\star_s)
  \;=\; \Big(\max_a D_0 - \E_{\pi_{\theta_0^\star}} D_0\Big)
     +\; \Big(\max_b D_1 - \E_{\pi_{\theta_1^\star}} D_1\Big)
     + O(s_0^4 + s_1^4) .
\]
Under (G2), with $\bar h_i, \bar w_i$ the heights and widths of the support
bumps of player $i$, each term equals
$\sum_k w^{(i)\star}_k h^{(i)}_k \big(1 - w^{(i)}_k /
\sqrt{(w^{(i)}_k)^2 + s_i^2}\big)
= \tfrac12 \bar h_i s_i^2 / \bar w_i^2 + O(s_i^4)$. In the symmetric two-peak
game with a common floor $s$ this collapses to the familiar
$2h\big(1 - w/\sqrt{w^2+s^2}\big)$.
\end{lemma}

\begin{proof}
Gaussian convolution: a bump $(h, w, p)$ convolved with $\mathcal{N}(0,s^2)$
has height $h w / \sqrt{w^2 + s^2} = h\,(1 - \tfrac{s^2}{2w^2} + O(s^4))$ and
retains its center; features that are polynomial pick up central moments, which
shift even-degree terms by $O(s^2)$ and odd-degree terms not at all (for a
general $C^2$ feature, $\E_{\mathcal{N}(x,s^2)}\varphi = \varphi(x) +
\tfrac{s^2}{2}\varphi'' + O(s^4)$, giving the same order). At $z^\star_s$ the
coupling contribution to player~0's best-response problem is
$\ip{\varphi(a)}{\E_{\pi_{\theta_1^\star}}\psi} = O(s_1^2)$ uniformly in $a$
(compactness plus boundedness of $\varphi$), so player~0's best response is
within $O(s_1^2)$ of a maximizer of $D_0$ and the loss against it is the
smoothing deficit of $D_0$ at $\mu^\star$'s atoms; the same computation with
$0 \leftrightarrow 1$ gives the second term. The two terms are computed from
different landscapes and different floors and are simply added --- there is no
factor $2$ except in the symmetric special case.
\end{proof}

Experiment 5 verifies the exact constant to $10^{-5}$ and the $s^2$ rate
(measured log--log slope $2.002$) in the symmetric instance.

\subsection{The local convergence theorem}

\begin{assumption}[Curvature dominance at the realized scale]\label{ass:curv}
There are radii $r_0 > 0$ and constants $\alpha_0, \alpha_1 > 0$ such that for
every configuration $z = (\theta_0, \theta_1)$ with each component mean within
$r_0$ of its assigned support point, each weight vector within $r_0$ of
$w^{(i)\star}$, and $\sigma^{(i)}$ at the floor $s_i$:
\[
  w^{(0)}_k\, \big(Q_{\nu}\big)_{s_0}''\big(\mu^{(0)}_k\big) \le -\alpha_0
  \quad \text{for all } k \le K_0,
  \qquad
  w^{(1)}_l\, \big(W_{\mu}\big)_{s_1}''\big(\mu^{(1)}_l\big) \le -\alpha_1
  \quad \text{for all } l \le K_1,
\]
where $\nu = \pi_{\theta_1}$, $\mu = \pi_{\theta_0}$ and $Q, W$ are the
effective landscapes \eqref{eq:effective}. (For multidimensional actions,
read $''$ as the Hessian and $\le -\alpha$ as ``$\preceq -\alpha I$''.)
\end{assumption}

Under (G1) the effective landscapes are \eqref{eq:effective-G1}, so
Assumption~\ref{ass:curv} reads, for player~0,
\[
  w^{(0)}_k\Big[ D_{0,s_0}''(\mu^{(0)}_k)
  + \big\langle \varphi_{s_0}''(\mu^{(0)}_k),\, \E_{\pi_{\theta_1}}
  \psi \big\rangle \Big] \le -\alpha_0 ,
\]
and for player~1 with $D_1$, $\psi$, and the moment $\E_{\pi_{\theta_0}}
\varphi$ entering with the opposite sign. This is where finite rank pays off:
the opponent's influence on player~0's curvature is entirely through the
$m$-vector $\E_\nu \psi$, whose norm is $O(s_1^2 + r_0)$ on the neighborhood by
Lemma~\ref{lem:bias} and continuity. Under (G2) the assumption holds with room
to spare: $|D_{0,s}''(p)| \approx \bar h_0 / \bar w_0^2$ ($\approx 156$ for
$\bar w_0 = 0.08$) while the coupling term is $O(c\, r_0 \|f''\|)$, so the
inequality holds whenever
\[
  r_0 \;\lesssim\; \min\Big\{ \sqrt{\bar w_0^2 + s_0^2},\quad
  \bar h_0 / \big(c\, \bar w_0^2 \|f''\|\big) \Big\},
\]
and symmetrically for player~1 with its own bumps. The first bound (the
concavity cap of the Gaussian bump) is the binding one for all moderate $c$;
the coupling can only bite once $c = \Omega(\bar h_0 / \bar w_0^2)$.
Experiment~3 measures exactly this: the certified region is flat in $c$ through
$c = 8$ and collapses only at extreme couplings. Note that in an asymmetric
game the two radii can differ by orders of magnitude --- the certified island is
the intersection, and the player with the flatter landscape is the binding one.

\begin{theorem}[Local monotonicity and linear convergence]\label{thm:local}
Under Assumptions~\ref{ass:reg} and~\ref{ass:curv}, let $N$ be the set of
configurations described there (a convex neighborhood of $z^\star_s$ in the
means and simplex coordinates). Then:
\begin{enumerate}
\item[(i)] The parametric field $F$, expressed in the mirror coordinates the
update actually uses (simplex weights for each categorical head, Fisher--Rao
metric for each set of means), is monotone on $N$; with magnet terms
$\tau_i > 0$ centered at any $\bar z \in N$, the regularized field
$F_\tau(z) = F(z) + \big(\tau_0 \nabla \KL_0,\ \tau_1 \nabla \KL_1\big)$ is
$\tau'$-strongly monotone on $N$ with
$\tau' \ge \min_i \tau_i\, \kappa_{N,i}$, where $\kappa_{N,i} > 0$ is the
strong-convexity modulus of player $i$'s KL regularizer on $N$.
\item[(ii)] The discrete update \eqref{eq:cat}--\eqref{eq:gauss} with magnets
fixed at $\bar z$ is a local contraction: for $\eta_i$ below a threshold set by
the Lipschitz constant of $F$ on $N$, the iterates converge linearly to the
unique fixed point $z_\tau(\bar z) \in N$ of $F_\tau$, at rate
$1 - \Theta(\min_i \eta_i \tau_i)$.
\item[(iii)] The magnet outer loop (refresh $\bar z \leftarrow z$ every $T$
steps) has $z^\star_s$ as its unique fixed point in $N$, and the composite
iteration initialized in a sufficiently small neighborhood of $z^\star_s$
converges to $z^\star_s$; by Lemma~\ref{lem:bias} the limiting policy pair is
an $O(s_0^2 + s_1^2)$-Nash equilibrium of the underlying continuous game.
\end{enumerate}
\end{theorem}

\begin{proof}
(i) By Lemma~\ref{lem:skewblock} it suffices to show that each player's payoff
is concave in its own coordinates on $N$, in the appropriate sign. We give
player~0's computation; player~1's is the same computation applied to $W_\mu$
in place of $Q_\nu$, with its own $K_1$, $s_1$, $\alpha_1$ --- \emph{not} an
appeal to symmetry.

Own coordinates are $(w^{(0)}, \mu^{(0)})$ with $\sigma^{(0)}$ pinned at the
floor (the update pushes $\rho^{(0)}$ below the floor throughout $N$ ---
widening a component strictly lowers the smoothed effective landscape near a
strict local max --- so the clip is active and the $\rho$-block is inert;
formally, monotonicity is needed only on the face of the box where the
dynamics live). On $N$, writing $\bar Q(x) \coloneqq (Q_\nu)_{s_0}(x)$,
\[
  \widehat U(w^{(0)}, \mu^{(0)}; \nu)
  = \sum_k w^{(0)}_k\, \bar Q\big(\mu^{(0)}_k\big)
  + \text{terms without } (w^{(0)}, \mu^{(0)}).
\]
The Hessian in $(w^{(0)}, \mu^{(0)})$ has blocks: $\partial^2_{ww} = 0$
(linear), $\partial^2_{\mu_k \mu_k} = w^{(0)}_k \bar Q''(\mu^{(0)}_k) \le
-\alpha_0$ by Assumption~\ref{ass:curv}, $\partial^2_{w_k \mu_k} =
\bar Q'(\mu^{(0)}_k)$, and zero between distinct components. Each
$2 \times 2$ block
$\begin{psmallmatrix} 0 & \bar Q' \\ \bar Q' & w \bar Q'' \end{psmallmatrix}$
is negative semidefinite iff $\bar Q' = 0$; on $N$,
$\bar Q'(\mu^{(0)}_k) = O(r_0 + s_1^2)$ rather than exactly zero, so we argue
in the mirror geometry instead. The update~\eqref{eq:cat} never uses curvature
of $\widehat U$ in $w^{(0)}$ --- it is the exact proximal step of the
\emph{linear} functional $\ip{w^{(0)}}{q^{(0)}}$ --- and the monotonicity gap
decomposes as
\[
  \ip{F(z) - F(z')}{z - z'} =
  \underbrace{\textstyle\sum_i \sum_k
  \big(w^{(i)}_k - w^{(i)\prime}_k\big)\big(q^{(i)\prime}_k - q^{(i)}_k\big)}
  _{\text{categorical, both players}}
  \;+\; \underbrace{\text{mean terms}}
  _{\ \ge\ \alpha_0 \|\mu^{(0)} - \mu^{(0)\prime}\|^2_{w}
    + \alpha_1 \|\mu^{(1)} - \mu^{(1)\prime}\|^2_{w}} .
\]
In the categorical sum, the part of $q^{(i)}$ that depends on the
\emph{opponent} is linear in the opponent's policy, and the two players'
contributions cancel in pairs by the computation of Lemma~\ref{lem:lift}
(Step~4--5 applied to the discretization induced by the two component sets ---
a $K_1 \times K_0$ rectangular matrix game, no squareness needed). The
remaining categorical--own-mean interaction is
$O(r_0)\,\|w^{(i)} - w^{(i)\prime}\|\,\|\mu^{(i)} - \mu^{(i)\prime}\|$ and is
dominated by $\alpha_i\|\mu^{(i)}-\mu^{(i)\prime}\|^2 +
\tau_i \kappa_{N,i}\|w^{(i)}-w^{(i)\prime}\|^2$ once $\tau_i > 0$ and
$r_0 \le \min_i \sqrt{\alpha_i \tau_i \kappa_{N,i}}$. Shrinking $r_0$
accordingly gives (i); the strong monotonicity constant is inherited from the
KL magnets, each $\kappa_{N,i}$-strongly convex on the truncated simplex and
Fisher--Rao ball, and the pair is $\min_i \tau_i \kappa_{N,i}$-strongly
monotone.

(ii) is the standard convergence of mirror descent / the forward step for
strongly monotone, Lipschitz variational inequalities (Sokota et al.\ 2023,
Thm.~3.4, applied on the convex set $N$; their proof is local and only needs
the VI restricted to $N$, plus the iterates staying in $N$, guaranteed for
small $\eta_i$ by the contraction toward $z_\tau(\bar z) \in N$). Distinct
per-player step sizes are absorbed by treating the product update as a
preconditioned forward step with block-diagonal preconditioner
$\mathrm{blockdiag}(\eta_0 I, \eta_1 I)$, whose condition number enters the
threshold; the rate is set by the slower player.

(iii) Fixed points of the outer loop satisfy $z = z_\tau(z)$, i.e.\ $z$ is the
$\tau$-regularized equilibrium with magnet at itself, which holds iff $F(z) =
0$ on $N$'s active face. On $N$ the only zero of $F$ is $z_s^\star$: for the
means, $\bar Q'$ has the sign of $p_k - \mu^{(0)}_k$ near each $p_k$ by
Assumption~\ref{ass:curv} (and likewise for player~1 near each $r_l$); for the
weights, stationarity of \eqref{eq:cat} forces the $q^{(i)}$ to be constant on
the support, which by Assumption~\ref{ass:reg}(ii) --- uniqueness of the fully
mixed equilibrium of the support-restricted matrix game $A$ --- pins
$w^{(i)} = w^{(i)\star}$ up to the $O(s^2)$ smoothing perturbation of $A$
(nonsingularity is an open condition, so the perturbed matrix game still has a
unique fully mixed equilibrium, at distance $O(s^2)$). Local convergence of the
composite follows from (ii) plus continuity of $\bar z \mapsto z_\tau(\bar z)$
with $\|z_\tau(\bar z) - z^\star\| \le \gamma\|\bar z - z^\star\|$,
$\gamma < 1$ (the magnet fixed-point map is a contraction toward the
equilibrium on monotone games).
\end{proof}

\begin{remark}[What the experiments certify]
Experiment~2 measures, at the numerically-obtained fixed point of the two-peak
game: own-mean Hessian eigenvalues $-49.64$ (matching $-w_k \bar h\,
\bar w/(\bar w^2+s^2)^{3/2}$ to four digits), cross term
$\bar Q'(\mu_k) = \mp 0.023 \approx 0$, one-step residual $10^{-4}$, and
spectral radius $\rho = 0.9913$ of the linearized update --- with
$1 - \rho \approx 0.0087 \approx \eta\tau/(1 + \eta\tau) = 0.0099$: the
contraction is exactly the magnet's strong monotonicity, as (ii) predicts. That
instance is symmetric, so the two players' constants coincide; the theorem
predicts that in an asymmetric instance the measured rate is set by
$\min_i \eta_i \tau_i$ and the certified radius by $\min_i r_{0,i}$.
\end{remark}

\subsection{Multiple islands: the same theorem explains the failures}

\begin{proposition}[Coexisting monotone islands]\label{prop:islands}
Let $\hat a \in \cA$ and $\hat b \in \cB$ and let $\hat z$ be the configuration
in which every component of player~0 sits at $\hat a$ and every component of
player~1 at $\hat b$ (any weights, $\sigma$ at the floors). Write
$\hat\mu = \delta_{\hat a}$-smoothed, $\hat\nu = \delta_{\hat b}$-smoothed for
the induced policies. Suppose
\begin{itemize}
\item[(a)] $\hat a$ is a strict local maximum of $(Q_{\hat\nu})_{s_0}$ and
$\hat b$ is a strict local maximum of $(W_{\hat\mu})_{s_1}$.
\end{itemize}
Then $\hat z$ is a fixed point of the parametric MMD dynamics, and
Assumption~\ref{ass:curv} holds in a neighborhood of $\hat z$ with
$\alpha_0 = |(Q_{\hat\nu})_{s_0}''(\hat a)| - O(r_0)$ and $\alpha_1$
analogously; hence $\hat z$ is locally monotone and locally exponentially
stable by the argument of Theorem~\ref{thm:local} --- \emph{regardless of
whether $\hat z$ realizes a Nash equilibrium}.
\end{proposition}

\begin{proof}
At $\hat z$ each player faces a fixed effective landscape ((a) is a condition
on that landscape, already including whatever the coupling contributes at the
opponent's current policy), so the mean gradients vanish, all $q^{(i)}_k$ are
equal within a player and the categorical step is stationary, and $\rho$ is
pinned at the floor: $\hat z$ is a fixed point. The curvature claim is (a), and
the neighborhood claims follow by continuity exactly as in
Theorem~\ref{thm:local}.
\end{proof}

\begin{remark}[The moment-matching special case]\label{rem:momentmatch}
Under (G1), a convenient sufficient condition for (a) is: $\E_{\hat\nu}\psi =
0$ and $\E_{\hat\mu}\varphi = 0$ (both couplings vanish identically), together
with $\hat a$ a strict local max of $D_{0,s_0}$ and $\hat b$ a strict local max
of $D_{1,s_1}$. This is how the codebase's traps arise, and it is the form in
which the phenomenon was first observed; but Proposition~\ref{prop:islands}
shows moment matching is \emph{not} necessary --- any pair of points that are
mutually stationary and curvature-dominated for the \emph{coupled} landscapes
generates an island. Experiment~8D found exactly such a non-moment-matched
trap (a mass-dominant decoy off the moment-matched point still captures the
dynamics, at an uglier configuration), which the symmetric moment-matching
statement could not account for.
\end{remark}

\begin{corollary}[The decoy trap]\label{cor:decoy}
In \texttt{DecoyWellGame} with peaks $(\pm 1)$ and uniform equilibrium weights,
the target first moment is $0 = \iota(0)$, so all mass at the decoy center
satisfies the condition of Remark~\ref{rem:momentmatch}; the decoy is a strict
local max of $D_s$ for every $s \ge 0$, giving (a). Hence ``everyone on the
decoy'' is a locally-monotone, locally-stable non-Nash fixed point with
exploitability $\approx (\max D_0 - D_0(0)) + (\max D_1 - D_1(0)) \approx 0.72$.
\end{corollary}

Experiment 4 confirms every claim numerically: at the trap endpoint the
moments vanish exactly, own-Hessians $-1.71 < 0$, spectral radius $0.9902$;
the Nash island of the same game coexists ($\rho = 0.9913$, Hessians
$-197.9$). The same experiment isolates the selection mechanism: from
identical means $\pm 1.5$, initial $\sigma = 1.5$ lands in the decoy island
(mass beats height under smoothing: the smoothed height of a bump decays as
its \emph{mass} $h w / s$) while initial $\sigma = 0.1$ lands in the Nash
island.

\begin{remark}[The other documented failures]
The remaining counterexamples in the codebase are boundary phenomena of the
same picture, not new mechanisms. \emph{Weight starvation}: the Fisher--Rao
mean-update \eqref{eq:gauss} carries the factor $w^{(i)}_k$ through
$\nabla_{\mu^{(i)}_k} \widehat U = w^{(i)}_k \bar Q'(\mu^{(i)}_k)$; as
$w^{(i)}_k \to 0$ the component's mean field vanishes --- the island containing
$z^\star$ requires all weights bounded away from $0$, and trajectories that hit
the simplex boundary leave every such island (the mirror geometry makes the
boundary repelling in $\log$-coordinates but the \emph{mean} freeze happens
first). This can afflict one player alone. \emph{Symmetric ties}:
configurations with exactly coincident components within a player are fixed by
an invariance of that player's own field (a measure-zero invariant set --- an
island of the degenerate kind); note this is a symmetry of the
\emph{parameterization}, present regardless of whether the game is symmetric.
\emph{Basin boundaries}: Experiment~2(D) shows the empirical basin of the Nash
island on the two-peak game is the product of the peaks' Voronoi cells
(predictor accuracy $96.7\%$, threshold insensitive on $[0.45, 0.55]$, i.e.\
half the peak separation) --- strictly larger than the certified monotone region
(the $\sqrt{\bar w^2 + s^2}$ concavity cap), as expected: monotonicity is
sufficient, not necessary, and outside the cap the funnel is maintained by the
(non-monotone but inward) effective-landscape pull.
\end{remark}

\begin{remark}[Where player-exchange symmetry \emph{would} buy
something]\label{rem:where-symmetry-helps}
The results above are local: they describe the islands, not which island a
given initialization reaches. The global (transit) analysis in the companion
notes reduces the two-player dynamics to $K$ independent single-component
scale-space flows on the smoothed landscape, and that reduction rests on the
coupling term vanishing \emph{along the whole trajectory}, not merely at the
equilibrium. In the codebase that holds because of a \emph{reflection}
symmetry --- even landscape, odd centered features, symmetrically placed
initial means --- which makes $\E_{\nu_t}\psi \equiv 0$ for all $t$; it is
strictly stronger than player-exchange symmetry and is broken by asymmetric
equilibrium weights (which is precisely the regime where the experiments find
the ``ejection'' failure mode that the flow reduction cannot see). Without it,
the correct general statement treats $\varepsilon(t) \coloneqq
\|\E_{\nu_t}\psi\|$ as an exogenous perturbation and requires a shadowing bound
of size $O(\varepsilon \sup\|\varphi'\| / \alpha)$ around the flow. That bound
is not proved here; it is the main open item. Nothing in
Sections~\ref{sec:layer1}--\ref{sec:local} depends on it.
\end{remark}

\begin{remark}[Separable multidimensional games]\label{rem:multidim}
For $u(a,b) = \sum_{d=1}^{n} u_d(a_d, b_d)$ with each $u_d$ of class
\eqref{eq:game} --- the per-axis landscapes and features may differ between the
players and between axes --- and policies that are $K_i$-mixtures of diagonal
Gaussians with a \emph{shared} categorical head, Lemmas~\ref{lem:lift}
and~\ref{lem:skewblock} apply verbatim, and Theorem~\ref{thm:local} holds at
any configuration whose component means sit on distinct support-point
\emph{vectors} (the own-mean Hessian is the direct sum of per-axis Hessians, so
Assumption~\ref{ass:curv} is the per-axis assumption plus a minimum over axes).
The shared head couples the axes only through $q^{(i)}_k = \sum_d q^{(i)}_{k,d}$,
which is why the codebase's Exp.~3 (\texttt{MultiDim.md}) sees weight
starvation \emph{amplified} by dimension --- a per-axis-factored head removes
the summation and restores the 1-D picture per axis. Two further quantitative
changes are recorded there: the smoothed amplitude of an isotropic bump decays
as $h\,(w/s)^{n}$, so the dominance condition of Assumption~\ref{ass:curv}'s
checkable form involves the effective mass $h w^{n}$; and the bias of
Lemma~\ref{lem:bias} gains a factor $n$ (trace over axes).
\end{remark}

\section{Summary of predictions and their experimental status}

\begin{enumerate}
\item[P1.] Tabular (measure-space) MMD converges globally on the very games
whose parametric versions have traps; skewness of the lifted field is exact.
\emph{Confirmed} (exp1: skewness $\sim 10^{-16}$; decoy well converges even
from the decoy; Forsaken converges under the step-size bound and cycles
outside it).
\item[P2.] The Nash-realizing configuration is a locally-monotone, locally
exponentially stable fixed point; contraction rate $\approx \eta\tau$;
own-Hessians match $-w_k D_s''(p_k)$. \emph{Confirmed} (exp2: $\rho = 0.9913$,
$1-\rho \approx \eta\tau/(1{+}\eta\tau)$; Hessians match to 4 digits).
\item[P3.] The decoy configuration is a second monotone island: mutually
stationary, concave own-blocks, $\rho < 1$, exploitability $0.72$.
\emph{Confirmed} (exp4), including the island-selection mechanism (initial
$\sigma$); exp8D additionally confirms the non-moment-matched variant
predicted by Proposition~\ref{prop:islands}.
\item[P4.] The certified island radius is $\min\{\sqrt{\bar w^2+s^2},\,
O(\bar h /(c \bar w^2))\}$ per player: flat in $c$ for moderate $c$, collapsing
only at extreme $c$. \emph{Confirmed} (exp3, three-peak game,
$\bar w = 0.08$): the certified radius is exactly $0.08 \approx
\sqrt{\bar w^2+s^2}$ and \emph{identical} for $c \in [0.5, 8]$, degrading only
at $c \in \{32, 100, 300\} = O(\bar h/\bar w^2) = O(156)$. The empirical basin
\emph{grows} with $c$ (the coupling force is the only inward pull outside the
well cap), and the total failure at $c = 8$, $\eta = 0.05$ is a \emph{step-size}
violation, not an island failure: the same inits converge on all seeds at
$\eta = 0.01$ (the Lipschitz constant grows with $c$, tightening
$\eta \lesssim \tau/L^2$).
\item[P5.] Residual exploitability at convergence is the sum of the two
players' smoothing biases, $\Theta(s_0^2) + \Theta(s_1^2)$. \emph{Confirmed}
(exp5, symmetric instance: matches to $10^{-5}$; log--log slope $2.002$).
\end{enumerate}

\paragraph{The one-sentence theory.} \emph{MMD converges on these games because
the games are monotone in the space of mixed strategies --- a consequence of the
zero-sum property alone, with no symmetry between the players --- and the
Gaussian-mixture policy is locally a faithful (Fisher--Rao) chart of that
space around any configuration at which each player's effective landscape is
curvature-dominated at its component means; the Nash is one such
configuration --- convergence to it is guaranteed exactly when the
initialization (its means \emph{and} its smoothing scale) selects the Nash
island rather than one of the finitely many other islands the same theorem
stabilizes.}

\begin{thebibliography}{9}
\bibitem{sokota2023} S.~Sokota, R.~D'Orazio, J.~Z.~Kolter, N.~Loizou,
M.~Lanctot, I.~Mitliagkas, N.~Brown, C.~Kroer.
\emph{A Unified Approach to Reinforcement Learning, Quantal Response
Equilibria, and Two-Player Zero-Sum Games.} ICLR 2023 (arXiv:2206.05825).
\bibitem{hsieh2021} Y.-P.~Hsieh, P.~Mertikopoulos, V.~Cevher.
\emph{The Limits of Min-Max Optimization Algorithms: Convergence to Spurious
Non-Critical Sets.} ICML 2021.
\bibitem{glicksberg1952} I.~L.~Glicksberg. \emph{A further generalization of
the Kakutani fixed point theorem, with application to Nash equilibrium
points.} Proc.\ AMS, 1952.
\bibitem{amari1998} S.~Amari. \emph{Natural gradient works efficiently in
learning.} Neural Computation, 1998.
\bibitem{hazan2016} E.~Hazan, K.~Y.~Levy, S.~Shalev-Shwartz. \emph{On
Graduated Optimization for Stochastic Non-Convex Problems.} ICML 2016
(arXiv:1503.03712).
\end{thebibliography}

\end{document}
